\section{实验装置与方法}

\subsection{测量线路}

实验采用以Model 128A型锁定放大器为核心的C--V测量线路（如图9-2-10所示）。测量线路主要包括待测pn结电容$C_{\mathrm{x}}$、已知标准电容$C_{0}$、直流偏置电路和保护电路四部分。交流测试信号$v(t)$由锁定放大器内部振荡器提供，频率$f_{\mathrm{s}} = 1133\,\mathrm{Hz}$，示波器测量峰-峰值$V_{\mathrm{so}} = 140.8\,\mathrm{mV_{pp}}$，锁定放大器测量有效值$V_{\mathrm{s}} = 49.47\,\mathrm{mV}$，满足小信号条件$v(t) \ll V_{\mathrm{D}} + V_{\mathrm{R}}$。锁定放大器参考相位基准$\phi_{0} = -171.93^{\circ}$。

待测pn结电容$C_{\mathrm{x}}$与已知电容$C_{0} = 5 \times 10^{3}\,\mathrm{pF}$串联后接入交流信号回路，满足$C_{0} \gg C_{\mathrm{x}}$的分压近似条件（待测电容零偏压下约几十至百余$\mathrm{pF}$）。$C_{0}$两端的交流电压$v_{\mathrm{i}}$送入锁定放大器信号输入端进行测量。由式(\ref{eq:meas-principle})，$v_{\mathrm{i}} \approx (C_{\mathrm{x}}/C_{0})\,v(t)$，理论灵敏度$v(t)/C_{0} = 9.894\,\mu\mathrm{V/pF}$。

直流反向偏压$V_{\mathrm{R}}$由晶体管直流稳压电源提供（$0$--$10\,\mathrm{V}$可调），通过$R_{2} = 100\,\mathrm{k}\Omega$隔离电阻加至pn结两端。隔离电阻的作用是防止交流信号被直流电源旁路，其阻值远大于pn结容抗$1/(\omega C_{\mathrm{x}})$。当pn结反向漏电流小于$1\,\mu\mathrm{A}$时，$R_{2}$上的直流压降可忽略。

锁定放大器输入端设有二极管保护电路：当输入信号超过约$300\,\mathrm{mV}$时，保护二极管导通，防止仪器因过载而损坏。

\subsection{锁定放大器}

锁定放大器是C--V测量的核心仪器。本实验使用Model 128A型锁定放大器，主要操作参数包括：

灵敏度（Sensitivity）选择需根据输入信号大小确定：在测量$v(t)$参考信号时，灵敏度置于$250\,\mathrm{mV}$档；测量$C_{0}$两端输出信号$v_{\mathrm{i}}$时，灵敏度需逐步提高至适当档位。时间常量（Time Constant）决定低通滤波器的等效噪声带宽$\Delta f_{\mathrm{N}} = 1/(4T)$，可根据信噪比要求在$1\,\mathrm{ms}$至$1000\,\mathrm{ms}$之间选取。参考通道内的移相器可在$0^{\circ}$--$360^{\circ}$范围连续调节，用于匹配输入信号与参考信号之间的相位差以获得最大直流输出。

本实验中锁定放大器输出两项正交分量：$R_{1\omega}$（与参考信号同相分量）和$R_{2\omega}$（正交分量）。当待测pn结无漏电时，电容性阻抗对应的输出信号与参考信号相位差约为$90^{\circ}$，此时$R_{1\omega} \propto C_{\mathrm{x}}$、$R_{2\omega} \approx 0$；若存在漏电，并联电阻分量将改变相位关系，$R_{2\omega}$不再为零。

\subsection{待测样品与标定}

实验待测样品为五个平面扩散工艺制作的$\mathrm{p}^+\mathrm{n}$单边突变结二极管（编号D1--D5），结面积$A = 5.03 \times 10^{-3}\,\mathrm{cm}^{2}$。在$V_{\mathrm{R}} = 10\,\mathrm{V}$反向偏压下，各二极管反向漏电流应小于$1\,\mu\mathrm{A}$。

C--V曲线的纵轴（电容轴）通过固定已知电容进行标定。在二极管处依次换接标称值为$20$、$40$、$60$、$80$、$100\,\mathrm{pF}$的标准电容，记录锁定放大器$R_{1\omega}$输出，通过线性拟合建立输出电压与电容值的换算关系$C_{\mathrm{x}} = (R_{1\omega} - b)/k$。此外，样品盒空挡的寄生电容$C_{\mathrm{stray}} = 1.4\,\mathrm{pF}$需从测量值中扣除。

\subsection{实验内容}

实验按照以下步骤进行：

\begin{enumerate}
  \item 了解Model 128A锁定放大器的工作原理与使用方法。用示波器测量交流测试信号$v(t)$的频率和峰-峰值，用锁定放大器测量其有效值$V_{\mathrm{s}}$和相位$\phi_{0}$。
  \item 观测锁定放大器从噪声中提取微弱信号的能力。将白噪声与$v(t)$叠加，在不同时间常量（$1$、$3$、$10$、$30$、$100$、$300$、$1000\,\mathrm{ms}$）下，用$X$--$Y$记录仪记录锁定放大器的$C$--$t$输出曲线，分析噪声标准差$\sigma$与等效噪声带宽$\Delta f_{\mathrm{N}}$的关系。
  \item 用固定已知电容（$20$--$100\,\mathrm{pF}$）检验锁定放大器$R_{1\omega}$输出与待测电容$C_{\mathrm{x}}$的线性关系，并标定电容轴。
  \item 以D2为典型样品，在不同频率（$933$--$1633\,\mathrm{Hz}$）下测量相位$\phi_{\mathrm{v}}$随反向偏压$V_{\mathrm{R}}$的变化，分析相位偏离$\phi_{0}$的原因。
  \item 对五个二极管逐一测量$R_{1\omega}$--$V_{\mathrm{R}}$曲线（$V_{\mathrm{R}}$由$0$至$10\,\mathrm{V}$，步进$0.2\,\mathrm{V}$），经标定和扣除寄生电容后得到各二极管的$C$--$V_{\mathrm{R}}$数据，进而计算杂质浓度分布$N(w)$。
\end{enumerate}
